\section{A uniform relative theorem in iPRA}
\label{sec:relative}
This section proves the new mathematical obligation directly from literal
rows. It does not assume any consistency statement or any theorem about the
register-program compiler. Let $P$ be either supplied witness-calculus parent,
$F$ its accelerated table, and $M$ the delivered marker-separated table.
Native start is 74 and native halt is 254 throughout. In particular, native
label 194 is an ordinary working control; it is not the halt label.

\subsection{The two accelerations}
The first transformation changes precisely
\[
 (26,1):1R42\longmapsto1R91,
 \qquad (101,0):0L26\longmapsto0R17.
\]
All other rows are copied. Here $wDq$ means write $w$, move in direction
$D$, and enter $q$. The unchanged old detours are
\[
\begin{array}{c|c}
\text{row}&\text{old action}\\\hline
(42,0)&0R13\\
(13,0)&0L91\\
(26,0)&0R194\\
(194,0)&0R17.
\end{array}
\]
We must justify the bits read on those detours, not just their control names.
For a configuration with head $h$, write $b_j$ for its bit at $h+j$.
The following three invariants suffice:
\begin{align*}
 q=101&\ \Longrightarrow\ b_{-1}=0,\tag{A}\\
 q=11&\ \Longrightarrow\ (b_{-1}=0\ \lor\ b_1=0),\tag{B}\\
 q=26&\ \Longrightarrow\
       \bigl(b_0=0\ \lor\ (b_1=0\land b_2=0)\bigr).\tag{C}
\end{align*}
They hold for both $P$ and $F$, starting in control 74 on any complete tape.
Indeed, the only incoming row to 101 is $(96,0):0R101$, proving (A).
Every incoming row to 11 writes zero: a right move gives the first
disjunct of (B), and a left move gives the second.
In $P$, the incoming rows to 26 are $(11,0):0L26$ and
$(101,0):0L26$; in $F$ only the former remains.
A 101 predecessor gives $b_0=0$ by (A).
At an 11 predecessor, the scanned zero becomes the new $b_1$; (B)
says either the new $b_0$ is zero or the new $b_2$ is zero.
This proves (C). Since 74 is none of these controls, the initial conditions
are vacuous. This is a one-step inductive argument over the actual rows.

The independent verifier also checks this induction step directly. It adds
both reader-history guards used below, examines every one of the 194
working controls and all $2^7$ radius-three binary windows, and tests
preservation whenever the guards hold. Current guards inspect radius two;
a single head move enlarges the needed window to radius three, so no
unexamined exterior bit affects the conclusion. For each parent and each
accelerated table this gives 24,832 examined state/window pairs, of which
24,560 satisfy the hypotheses, with every such successor satisfying the
guards. The receipt is \path{verification/independent.json}. This finite
certificate checks the one-step implication used by the induction; it
is not an extrapolation from a sampled execution.

At $(26,1)$, (C) provides two zeros immediately right of the head.
The parent takes the three moves $R,R,L$ through 42 and 13; the child
moves directly right into 91. Both tapes and heads agree exactly.
At $(101,0)$, (A) provides a zero to the left. The parent takes $L,R,R$
through 26 and 194; the child moves directly right into 17. Again the
complete tape agrees. Neither three-step segment visits halt.
The independent checker enumerates the 12 compatible radius-two windows,
checking all cells and every read location. Arbitrary exterior tape is
untouched. The invariants, rather than this window enumeration alone,
establish that these cases cover executions.

For the arithmetic formalization, encode each finite-support tape by its
canonical left and right binary integers relative to the head, together
with the control and an integer head coordinate. Blank suffixes are
implicit. For a fixed head position, two coded tapes agree at every cell
if and only if their two binary integers agree. Configuration equality and the update of these
codes are primitive recursive. The arbitrary-exterior local arguments
above specialize to these encodings, including nonstandard coded values.

Totalize halting by an absorbing configuration, without adding a working
state to the literal table. Let $C_T(n)$ denote the configuration of table
$T$ after $n$ totalized steps. Define a primitive-recursive duration $a(c)$
to be 3 at either accelerated row and 1 otherwise, including at halt, and
\[
 s(0)=0,\qquad s(n+1)=s(n)+a(C_F(n)).
\]
The preceding local identities prove
\[
 C_P(s(n))=C_F(n),\qquad n\le s(n)\le3n.\tag{1}
\]
The induction predicate consists of primitive-recursive equalities,
inequalities, and finite Boolean conditions. It is therefore available in
iPRA; if a chosen presentation permits only equations in its induction
schema, encode the displayed finite conjunctions and disjunctions by their
primitive-recursive Boolean characteristic functions.

\subsection{Separating the marker from the history selector}
Define the control map
\[
 f(51)=400,\qquad f(301)=88,\qquad
 f(q)=q\quad(q\notin\{51,301\}).
\]
The parent does not contain 400. Copy every row of $F$ except the two rows
at 301, applying $f$ to its source and destination controls. Then install
\[
\begin{array}{c|cc}
 &\text{read 0}&\text{read 1}\\\hline
88&0L51&1R103\\
51&0R122&1R51\\
400&0R16&1R400.
\end{array}\tag{2}
\]
Control 400 now performs the original register-marking scan. Control 51
is a dedicated history selector. All genuine old marker-entry edges point
to 400. Thus the erroneous blank-initialization behavior of the older
193-state core is not inherited.

Exactly one old working control, 301, is removed and one new control,
400, is added. The literal count remains 194. Every one of the 384 rows
of $F$ outside controls 301 and 88 commutes with $f$ in one step.
The remaining rows of $F$ needed for the four exceptional cases are
\[
\begin{array}{c|cc}
 &0&1\\\hline
301&0R45&1R103\\
88&0R122&1R103\\
122&0L137&0L1\\
137&0R30&1R70\\
30&0R45&1R46.
\end{array}
\]
Every incoming edge to 301 writes zero and moves right; every incoming edge
to 88 writes one and moves right. Neither is the initial control.
Consequently, at a parent occurrence of 301 the preceding bit is zero,
and at an occurrence of 88 it is one. This invariant is valid for every
complete initial tape, independently of source reachability.

Put the current head at coordinate zero. The four cases are
\[
\begin{array}{cccrl}
\text{old control}&b_{-1}&b_0&\text{child steps}&\text{common exit}\\\hline
301&0&0&5&(45,1)\\
301&0&1&1&(103,1)\\
88&1&0&3&(122,1)\\
88&1&1&1&(103,1).
\end{array}\tag{3}
\]
Both bits are preserved; all other cells are untouched. For the two longer
cases, the control/head trajectories are
\begin{align*}
(88,0)&\to(51,-1)\to(122,0)\to(137,-1)
             \to(30,0)\to(45,1),\\
(88,0)&\to(51,-1)\to(51,0)\to(122,1).
\end{align*}
Every read is at $-1$ or zero. No intermediate state is halt. This explicit
check against the final table also rules out an interaction between the
simultaneously installed changes.

Extend $f$ to a configuration map $\widehat f$ by leaving tape and head
unchanged. On a configuration of $F$, let $r$ be 5 at $(301,0)$, 3 at
$(88,0)$, and 1 otherwise, including at halt. Set
\[
 u(0)=0,\qquad u(n+1)=u(n)+r(C_F(n)).
\]
The incoming-history invariant and (3) give the iPRA induction
\[
 C_M(u(n))=\widehat f(C_F(n)),\qquad
 n\le u(n)\le5n.\tag{4}
\]
The map $f$ sends no working control to halt. Its noninjectivity at
301 and 88 causes no difficulty: the old run and its incoming-history
invariant determine the appropriate local case.

\subsection{Bounds at every time, including intermediate child times}
Let $H_T(t)$ mean that $T$ has halted by time $t$ in its totalized blank
execution. Totalization makes $H_T$ monotone. From (1),
\[
 H_P(t)\to H_F(t),\qquad H_F(t)\to H_P(3t).
\]
For the first implication, $t\le s(t)$ gives $H_P(s(t))$, and then (1)
gives $H_F(t)$. For the second, use (1) and $s(t)\le3t$.
Similarly, (4) gives
\[
 H_F(t)\to H_M(5t),\qquad H_M(t)\to H_F(t).
\]
In particular, the reverse implication covers a halt at an arbitrary
intermediate child time. It does not assume that the alleged halt is a
synchronization boundary: $t\le u(t)$ and absorption transport it to one.
No inverse of a clock or unbounded least-time search is needed.

\begin{theorem}[New bounded witness transport]
\label{thm:relative}
For each exact parent-child pair in this archive,
\[
 \boxed{\quad
 \ipra\vdash H_P(t)\to H_M(5t),\qquad
 \ipra\vdash H_M(t)\to H_P(3t).
 \quad}
\]
Consequently
\[
 \ipra\vdash\NH(P)\leftrightarrow\NH(M),
 \qquad\NH(T):=\forall t\,\neg H_T(t).
\]
\end{theorem}
\begin{proof}
Compose the four bounded implications. To obtain either nonhalting
implication, assume a halting witness for the other machine and apply its
explicit witness bound. This is intuitionistically valid implication
reasoning. The construction and proof use primitive recursion and
quantifier-free induction; they use neither reflection for PA nor an
assumption that PA is consistent.
\end{proof}
